% PAGES: 46
% Continues section 1.5 "Lower triangular and upper triangular matrices.
% Tridiagonal matrices." begun in sec01_5.tex (ends there with the banded
% matrix definition and a band-$2$ example). Do not open sec01_5.tex.

A matrix of band $1$ is called a tridiagonal matrix and appears frequently
in practical applications.

\begin{definition}
The $n\times n$ matrix $A$ is tridiagonal if and only if
\[
A(j,k) \;=\; 0, \quad \forall\, 1\le j,k\le n \text{ with } |j-k|>1,
\]
or, equivalently, if and only if the only nonzero entries of $A$ could be
the main diagonal entries $A(i,i)$, for $i=1:n$, the upper diagonal
entries $A(i,i+1)$, for $i=1:(n-1)$, and the lower diagonal entries
$A(i,i-1)$, for $i=2:n$.
\end{definition}

\begin{examplex}
The positions of the entries of a $6\times 6$ tridiagonal matrix which
could be nonzero are denoted by ``$\times$'' in the matrix below:
\[
\begin{pmatrix}
\times & \times & 0 & 0 & 0 & 0 \\
\times & \times & \times & 0 & 0 & 0 \\
0 & \times & \times & \times & 0 & 0 \\
0 & 0 & \times & \times & \times & 0 \\
0 & 0 & 0 & \times & \times & \times \\
0 & 0 & 0 & 0 & \times & \times
\end{pmatrix}.
\]
\end{examplex}

\section{References}

Several linear algebra books are close in spirit to our approach
emphasizing numerical implementation and applications: Gilbert Strang's
``Introduction to Linear Algebra'' \cite{41} and ``Linear Algebra and Its
Applications'' \cite{40} and Peter Lax's ``Linear Algebra and Its
Applications'' \cite{26} are three acclaimed such texts; also, Strang's
OpenCourseWare MIT videos are a remarkable resource. A classical approach
to graduate level linear algebra can be found in Horn and Johnson
\cite{22}.

The numerical analysis approach to linear algebra was pioneered by the
classical monograph of Golub and Van Loan \cite{18}, and further
solidified by Demmel \cite{13}. The insightful numerical linear algebra
textbook by Trefethen and Bau \cite{43} presents the intuition behind
numerical linear algebra methods.

Some of the exercises contained in this chapter and throughout the book
are frequently asked in interviews for positions requiring quantitative
skills; see Stefanica, Radoičić, and Wang \cite{39} for more
interview--style questions.
